% Checked transcription from Coates--Kasprzyk--Pitton--Tveiten, published Proc.A477(2021)20210584, CC BY4.0.
% arXivv1 native TeX was a PRIVATE transcription aid, not the licensed export source.
\begingroup
\section*{Separately attributed earlier human proof component}
\noindent\textbf{Authors:} Tom Coates, Alexander M. Kasprzyk, Giuseppe Pitton and Ketil Tveiten.\par
\noindent\textbf{Work/version:} \emph{Maximally mutable Laurent polynomials}, published \emph{Proceedings of the Royal Society A}477(2021), article20210584, DOI\url{https://doi.org/10.1098/rspa.2021.0584}.\par
\noindent\textbf{Published source:} \url{https://wrap.warwick.ac.uk/196397/1/rspa.2021.0584.pdf}. Physical and printed pages coincide. Necessary definitions3,6,9--12; Proposition3.7 statement12; complete P-height/base/induction proof13--15. Published first page grants CC BY4.0.\par
\noindent\textbf{Editorial scope:} This is checked editable transcription of that published passage, with original mathematical content, notation, equations and conventions. A preprint TeX download was a local transcription aid only; its separate nonexclusive licence is not exported or relabelled CC. Prior normalization, non-negative-integer coefficient and monomial-span conventions are retained below. Lutz's current complex/non-normalized and zero-boundary adaptation remains literally unchanged in the preceding author block. This component supplies the cited coefficient-elimination mechanism; no editor hypothesis-transfer argument or proof is inserted. Optional Example3.8 and later Proposition3.9/classification are outside the selected support.\par
\medskip\hrule\medskip
\setcounter{equation}{0}
\renewcommand{\theequation}{3.\arabic{equation}}

% Published component: lattice/Fano definition; physical pages 2,3
We now attempt to reverse the construction described above. Let $N$ be a lattice of rank $n$, and let $P\subset{N_{\mathbb{Q}}}:=N\otimes_\mathbb{Z}\mathbb{Q}$ be a convex lattice polytope.
\par\noindent\textbf{Definition 1.2.} 
A convex lattice polytope $P\subset{N_{\mathbb{Q}}}$ is called \emph{Fano} if it is of maximum dimension with respect to the ambient lattice $N$, if it contains the origin $0$ in its interior ${P}^{\circ}$, and if the vertices of $P$ are primitive lattice elements. 

\par

% Published component: mutable-pair definition; physical pages 6
\par\noindent\textbf{Definition 1.7.} 
Let $N$ be a lattice and let $w\in M$ be a primitive vector in the dual lattice. Then $w$ induces a grading on $\mathbb{C}[N]$. Let $a\in\mathbb{C}[w^\perp\cap N]$ be a Laurent polynomial in the zeroth piece of $\mathbb{C}[N]$, where $w^\perp\cap N=\{v\in N\mid w(v)=0\}$. The pair $(w,a)$ defines an automorphism of $\mathbb{C}(N)$ given by
\[
\mu_{w,a}:\mathbb{C}(N)\rightarrow\mathbb{C}(N),\qquad\bm{x}^v\mapsto\bm{x}^v a^{w(v)}.
\]
Let $f\in\mathbb{C}[N]$. We say that $f$ is \emph{mutable with respect to $(w,a)$} if
\[
g:=\mu_{w,a}(f)\in\mathbb{C}[N],
\]
in which case we call $g$ a \emph{mutation of $f$} and $a$ a \emph{factor}.

\par

% Published component: normalization and original coefficient conventions; physical pages 9
Before introducing the mutation graph, we fix normalization conditions for the Laurent polynomials that we consider.

\par\noindent\textbf{Definition 2.2.}  \label{def:normalization}
A Laurent polynomial $f\in\mathbb{C}[N]$ is \emph{normalized} if for all vertices $v$ of $\operatorname{Newt}(f)$, the coefficient of the monomial $\bm{x}^v$ in $f$ is~$1$.

\par

\par\noindent\textbf{Convention 2.3.} 
From here onwards, we assume that all Laurent polynomials (and all mutation factors) are normalized. Similarly, although our Laurent polynomials are defined over $\mathbb{C}$, our expectation is that Laurent polynomials that are mirror to Fano manifolds have coefficients that are non-negative integers. From here onwards, we will require that all Laurent polynomials (and all mutation factors) have non-negative integer coefficients. Recall our standing convention that if $f \in \mathbb{C}[N]$ then the exponents of monomials in $f$ generate $N$.

\par

% Published component: cone/singularity-content terminology; physical pages 10,11
An edge $e$ of a Fano polygon $P$ determines the cone $\sigma$ over $e$, and hence a torus-fixed point in the toric variety $X_P$. A neighbourhood of this point is the affine toric variety $X_\sigma$ determined by $\sigma$, that is, the cyclic quotient singularity $\frac{1}{b}(1,c-1)$ where the rays $\{u,v\}$ of $\sigma$ are such that $u$,~$v$, and $\frac{c-1}{b} u + \frac{1}{b}v$ generate $N$. Write $b = dr$ and $c = dk$ where $d = \gcd\{b,c\}$. Then $X_\sigma$ is the cyclic quotient singularity $\frac{1}{dr}(1,kd-1)$. Furthermore, $r$ is the lattice height of $e$ above the origin, and $d$ is the lattice length of $e$.

When $d=nr$ for some $n\in\mathbb{Z}$, the singularity $X_\sigma$ is called a \emph{T-singularity} and $\sigma$ is called a \emph{$T$-cone}. If $n=1$ then $X_\sigma$ is a \emph{primitive} $T$-singularity and $\sigma$ is a \emph{primitive} $T$-cone. $T$-singularities are qG-smoothable~\textup{[28]}. A general $T$-singularity $\frac{1}{n r^2}(1, knr-1)$ admits a crepant partial resolution into $n$ primitive $T$-singularities $\frac{1}{r^2}(1, kr-1)$; this amounts to the fact that, since $d = n r$, the cone $\sigma$ can be decomposed as a union of $n$ primitive $T$-cones -- see figure 2.
When $d = s$ with $0 < s < r$, the singularity $X_\sigma$ is called an \emph{R-singularity} and $\sigma$ is called an $R$-cone. $R$-singularities are qG-rigid. When $d = nr+s$ for some non-negative $n \in \mathbb{Z}$ and $0 < s < r$, the singularity $X_\sigma = \frac{1}{dr}(1, kd-1)$ admits a (non-unique) crepant partial resolution into $n$ primitive $T$-singularities and the single $R$-singularity $\frac{1}{sr}(1,ks-1)$; this corresponds to the (non-unique) decomposition of $\sigma$ into $n$ primitive $T$-cones and one $R$-cone -- see figure 3.
\noindent The singularity $X_\sigma$ qG-deforms to the R-singularity $\frac{1}{sr}(1,ks-1)$. We define the \emph{residue} of~$\sigma$ to be
\[
\operatorname{res}(\sigma)=\begin{cases}
\emptyset&\text{if $s=0$}\\
\frac{1}{sr}(1,ks-1)&\text{otherwise,}
\end{cases}
\]
and the \emph{singularity content} of~$\sigma$ to be the pair $(n,\operatorname{res}(\sigma))$.


\par\noindent\textbf{Definition 3.2.} 
Let $P$ be a Fano polygon with edges $e_1,\ldots,e_m$, and suppose that the singularity content of the cone over $e_i$ is $(n_i, S_i)$. Let $n = n_1 + \cdots + n_m$, and let $\mathcal{B}$ be the multiset containing those $S_i$ such that $S_i \ne \emptyset$, counted with multiplicity. The \emph{singularity content} of $P$ is
\[
\operatorname{SC}(P)=(n, \mathcal{B}).
\]

\par

% Published component: residual definitions; physical pages 12
\par\noindent\textbf{Definition 3.5.} 
Let $C$ be an $R$-cone with primitive ray generators $\rho_1$,~$\rho_2 \in N$. Let $u \in M$ be the primitive normal vector to the edge $\operatorname{conv}\{\rho_1,\rho_2\}$ that is positive on $C$, and let $r = u(\rho_1) = u(\rho_2)$. The set $\mathcal{R}_C$ of \emph{residual points} of $C$ is
\[
\mathcal{R}_C := C \cap \{p \in N \, \mid \, u(p) \leq r\} \setminus \{0, \rho_1, \rho_2\}.
\]

\par

\par\noindent\textbf{Definition 3.6.} 
Let $P \subset {N_{\mathbb{Q}}}$ be a Fano polygon with singularity content $(n, \mathcal{B})$, and let $k = |\mathcal{B}|$. We say that $\mathcal{R} \subset N$ is \emph{a choice of residual points} for $P$ if and only if there exists a crepant subdivision of the spanning fan for $P$ into $n$ $T$-cones and $k$ $R$-cones $C_1,\ldots,C_k$ such that 
\[
\mathcal{R} = \mathcal{R}_{C_1} \cup \cdots \cup \mathcal{R}_{C_k}.
\] 

\par
\noindent Any two choices $\mathcal{R}_1$,~$\mathcal{R}_2$ of residual points for $P$ satisfy $|\mathcal{R}_1| = |\mathcal{R}_2|$.

% Published component: Proposition3.7 statement; physical pages 12
\begin{CoatesProp} 
Let $P$ be a Fano polygon. For each edge $e$ of $P$, fix a primitive vector $v_e$ in the direction of $e$ and write $w_e$ for the primitive inward-pointing normal vector to $e$. Write $(n_e, S_e)$ for the singularity content of the cone over $e$, and set $a_e = (1+\bm{x}^{v_e})^{n_e}$. There exists a Laurent polynomial~$f$ with Newton polytope $P$ and zero constant term that is mutable with respect to $(w_e, a_e)$ for each edge~$e$. Furthermore, the free parameters in the coefficients of the general such $f$ are in bijection with any choice of residual points of~$P$, and are affine-linear functions of the coefficients of such points.
\end{CoatesProp}
\begin{center}
\includegraphics[width=.55\linewidth]{assets/coates-figure5.pdf}\par
\textbf{Published Figure5.} Lattice points on $L$ outside the cone $C_e$ are at $P$-height greater than $H$.\par
\textit{Literal printed diagram label:} $L$.
\end{center}


% Published component: full P-height/base/induction proof; physical pages 13,14,15
\begin{proof}[Proof of Proposition~\ref{coates:prop:local_MMLP}]
Let us define the \emph{$P$-height} of a lattice point $p \in P$ to be the non-negative rational number $r$ such that $p$ lies on the boundary of $rP$. The set of $P$-heights of points in $P \cap N$ is a finite subset of $[0,1]$. Fix a set of residual points $\mathcal{R}$ for $P$, and write $\mathcal{R}_{\geq h}$ for the set of points in $\mathcal{R}$ of $P$-height at least $h$. Set $f = \sum_{v \in P \cap N} a_v \bm{x}^v$. We will prove the following statement by descending induction on $h$:
  \begin{equation}
  \tag{$\ddagger$}
  \label{coates:induction statement}
  \begin{minipage}{0.8\textwidth}
    The coefficients $a_v$ such that $v$ has $P$-height $h$ are affine-linear functions of the coefficients $a_w$ for $w \in \mathcal{R}_{\geq h}$, and these coefficients $a_w$ are independent.
  \end{minipage}
  \end{equation}
  \noindent To prove the Proposition, we need to prove statement \eqref{coates:induction statement} for all $h \geq 0$. It holds trivially for $h=0$.

  \subsubsection*{(ii) Base case: $h=1$} Points of $P$-height $1$ lie on the boundary of $P$. Consider an edge $e$ of $P$. Without loss of generality we may assume that $e$ is horizontal and at positive vertical height $r$ above the origin. Up to overall multiplication by a monomial, the coefficients of $f$ along $e$ are
  \begin{equation}
    \label{coates:eq:edge coefficients}
    1 + b_1 x + \cdots + b_{d-1} x^{d-1} + x^d  
  \end{equation}
  where $d = n_e r+s$ and $0 \leq s < r$, for some $b_1,\ldots,b_{d-1}$. The mutability condition gives that 
  \begin{equation}
    \label{coates:eq:edge equation}
    1 + b_1 x + \cdots + b_{d-1} x^{d-1} + x^d = 
    (1+x)^{r n_e} \sum_{i=0}^s c_i x^i
  \end{equation}
  for some coefficients $c_i$. This forces $c_0 = c_s = 1$. If $s = 0$ or $s = 1$ then there are no residual points on $e$, and \eqref{coates:induction statement} holds as \eqref{coates:eq:edge coefficients} coincides with $(1+x)^d$. If $s>1$ then to prove \eqref{coates:induction statement} it suffices to show that we can solve \eqref{coates:eq:edge equation} uniquely for $c_i$ in terms of the coefficients $b_{k_1}, \ldots, b_{k_{s-1}}$ of residual points on $e$. That is, it suffices to show that the $(s-1) \times (s-1)$ matrix with $(i,j)$ entry
  \[
      { n_e r \choose k_i - j}
  \]
  is invertible. But the determinant of this matrix is positive -- indeed it counts certain semi-standard Young tableaux, as can be seen by specializing to $1$ all variables in the Giambelli formula~\textup{([29], eqn A6)} that expresses Schur polynomials in terms of elementary symmetric polynomials. Thus statement \eqref{coates:induction statement} holds for $h=1$. 

  \subsubsection*{(iii) Induction step} Suppose that $H$ satisfies $0 < H < 1$ and that statement \eqref{coates:induction statement} holds for all $h > H$. We will prove \eqref{coates:induction statement} for $h=H$. Let $v \in P$ be a lattice point of $P$-height $H$. Since $P$ is Fano, $v$ lies in the cone $C_e$ over a unique edge $e$. Without loss of generality we may assume that $e$ is horizontal and at positive vertical height $r$ above the origin. Let $d$ denote the lattice length of $e$, and write $d = n_e r + s$ with $0 \leq s < r$. Consider the horizontal line $L$ through $v$; this is at vertical height $Hr$ above the origin. Lattice points on $L$ fall into three classes: residual points in $C_e$, non-residual points in $C_e$, and points outside $C_e$. Convexity implies that the points on $L$ outside $C_e$ are at $P$-height greater than $H$: see Figure~\ref{coates:fig5}. By the induction hypothesis, therefore, coefficients of these points are fixed as affine-linear combinations of coefficients of residual points of height greater than $H$. 
  
  Up to overall multiplication by a monomial, the coefficients of $f$ along $L$ are
  \[
    b_0 + b_1 x + \cdots + b_m x^m  
  \]
  for some $m$ and $b_0,\ldots,b_m$. The mutability condition gives that 
  \begin{equation}
    \label{coates:eq:L equation}
    b_0 + b_1 x + \cdots + b_m x^m = 
    (1+x)^{r H n_e } \sum_{i=0}^s c_i x^i
  \end{equation}
  for some coefficients $c_i$. A primitive $T$-cone with lattice height $r$ contains exactly $h$ lattice points at lattice height $h$, if $0 < h < r$, and so if $L \cap P$ contains no residual points and no points outside $C_e$ then $m<r H n_e$. In this case \eqref{coates:eq:L equation} gives that $b_0 = \cdots = b_m = 0$, and statement \eqref{coates:induction statement} holds. Otherwise \eqref{coates:eq:L equation} holds with $s \geq 0$ and $s+1$ equal to the total number of points in $L \cap P$ that are either residual or lie outside $C_e$. Statement \eqref{coates:induction statement} for $h=H$ follows if we can solve \eqref{coates:eq:L equation} uniquely for $c_i$ in terms of the coefficients $b_{k_0}, \ldots, b_{k_s}$ of points on $L$ that are either residual or lie outside $C_e$. For this, we use the same matrix-invertibility argument as above. This completes the induction step, and the proof of Proposition~\ref{coates:prop:local_MMLP}.
\end{proof}
\medskip\noindent\textbf{Earlier published reference[29]:} Fulton W, Harris J.1991 \emph{Representation theory}. In \emph{Graduate texts in mathematics}, vol.129. A First Course, Readings in Mathematics. New York, NY: Springer-Verlag. The author's exact Giambelli locator is eqn A6.\par
\noindent\textbf{Earlier published reference[28]:} Koll\'ar J, Shepherd-Barron NI.1988 Threefolds and deformations of surface singularities. \emph{Invent.Math.}91,299--338. DOI\url{https://doi.org/10.1007/BF01389370}.\par
\endgroup
